
%(BEGIN_QUESTION)
% Copyright 2006, Tony R. Kuphaldt, released under the Creative Commons Attribution License (v 1.0)
% This means you may do almost anything with this work of mine, so long as you give me proper credit

Når hydrogensulfid (H$_{2}$S) brennes med oksygen (O$_{2}$) under ideelle forhold, blir produktene vanndamp (H$_{2}$O) og svoveldioksid (SO$_{2}$). Skriv en balansert reaksjonsligning med alle reaktanter og produkter i riktige proporsjoner.

\vskip 10pt

Bestem også hvor mange mol oksygengass som trengs for fullstendig forbrenning av 15 mol hydrogensulfid.

\vskip 10pt

\underbar{file i00898}
%(END_QUESTION)





%(BEGIN_ANSWER)

22.5 mol oksygen (O$_{2}$) trengs for fullstendig forbrenning av 15 mol hydrogensulfid (H$_{2}$S).

%(END_ANSWER)





%(BEGIN_NOTES)

Balansering av reaksjonen med simultane lineære ligninger:

% No blank lines allowed between lines of an \halign structure!
% I use comments (%) instead, so Tex doesn't choke.

$$\vbox{\offinterlineskip
\halign{\strut
\vrule \quad\hfil # \ \hfil & 
\vrule \quad\hfil # \ \hfil & 
\vrule \quad\hfil # \ \hfil & 
\vrule \quad\hfil # \ \hfil & 
\vrule \quad\hfil # \ \hfil \vrule \cr
\noalign{\hrule}
%
% First row
1 & x & = & $y$ & $z$ \cr
%
\noalign{\hrule}
%
% Another row
H$_{2}$S & O$_{2}$ & $\to$ & H$_{2}$O & SO$_{2}$ \cr
%
\noalign{\hrule}
} % End of \halign 
}$$ % End of \vbox

% No blank lines allowed between lines of an \halign structure!
% I use comments (%) instead, so Tex doesn't choke.

$$\vbox{\offinterlineskip
\halign{\strut
\vrule \quad\hfil # \ \hfil & 
\vrule \quad\hfil # \ \hfil \vrule \cr
\noalign{\hrule}
%
% First row
{\bf Element} & {\bf Balanseligning} \cr
%
\noalign{\hrule}
%
% Another row
Hydrogen & $2 + 0x = 2y + 0z$ \cr
%
\noalign{\hrule}
%
% Another row
Oksygen & $0 + 2x = y + 2z$ \cr
%
\noalign{\hrule}
%
% Another row
Svovel & $1 + 0x = 0y + 1z$ \cr
%
\noalign{\hrule}
} % End of \halign 
}$$ % End of \vbox

Fra hydrogenbalansen ($2 + 0x = 2y + 0z$) får vi $y = 1$. Fra svovelbalansen ($1 + 0x = 0y + 1z$) får vi $z = 1$. Setter vi dette inn i oksygenbalansen ($0 + 2x = (1) + (2)(1)$), får vi $x = {3 \over 2}$. For heltallskoeffisienter uttrykker vi løsningen som $x = 3$, $y = 2$ og $z = 2$, med opprinnelig konstant for hydrogensulfid økt fra 1 til 2:

$$\hbox{2H}_2\hbox{S} + \hbox{3O}_2 \to \hbox{2H}_2\hbox{O} + \hbox{2SO}_2$$

Fra den balanserte ligningen ser vi molforholdet H$_{2}$S:O$_{2}$ = 2:3. Oksygenbehov for 15 mol hydrogensulfid beregnes slik:

$$\left({15 \hbox{ mol H}_2\hbox{S} \over 1}\right) \left({3 \hbox{ mol O}_2 \over 2 \hbox{ mol H}_2\hbox{S}}\right) = 22.5 \hbox{ mol O}_2$$









\vfil \eject

\noindent
{\bf Forberedende quiz:}

Velg den {\it balanserte} reaksjonsligningen for forbrenning av hydrogensulfid:

\begin{itemize}
\item{} $\hbox{3H}_2\hbox{S} + \hbox{O}_2 \to \hbox{2H}_2\hbox{O} + \hbox{3SO}_2$
\vskip 10pt
\item{} $\hbox{2H}_2\hbox{S} + \hbox{3O}_2 \to \hbox{2H}_2\hbox{O} + \hbox{2SO}_2$
\vskip 10pt
\item{} $\hbox{H}_2\hbox{S} + \hbox{O}_2 \to \hbox{2H}_2\hbox{O} + \hbox{SO}_2$
\vskip 10pt
\item{} $\hbox{2H}_2\hbox{S} + \hbox{O}_2 \to \hbox{3H}_2\hbox{O} + \hbox{3SO}_2$
\vskip 10pt
\item{} $\hbox{3H}_2\hbox{S} + \hbox{3O}_2 \to \hbox{2H}_2\hbox{O} + \hbox{3SO}_2$
\vskip 10pt
\item{} $\hbox{H}_2\hbox{S} + \hbox{2O}_2 \to \hbox{H}_2\hbox{O} + \hbox{2SO}_2$
\end{itemize}


%INDEX% Chemistry, stoichiometry: balancing a chemical equation

%(END_NOTES)
